\section{Introduction}
Daily athlete-monitoring systems commonly capture internal training load and self-reported wellness to support training planning and athlete communication \citep{Halson2014,Brink2010,Rago2020,Saw2016,Thorpe2015}. Training duration, rating of perceived exertion (RPE), session-RPE load, sleep, fatigue, soreness, mood, stress, and perceived readiness can be collected repeatedly in field settings. However, repeated self-reporting can impose reporting and data-management demands and frequently produces missing observations, requiring that monitoring studies define the intended decision problem before assuming that a larger panel is necessarily more useful \citep{Saw2016,Midoglu2024}.

A central practical question is therefore not simply whether workload and wellness correlate with readiness. Several studies indicate that wellness responses can vary with training exposure and that dose--response relations may be nonlinear or context-dependent \citep{Campbell2020,Campbell2021}. Rather, the relevant question for day-to-day monitoring is whether a larger panel yields an incrementally better forecast than a strong comparator that already uses the athlete's own recent readiness. If an athlete's preceding readiness reports already capture much of the persistent and short-term information relevant to tomorrow's report, a complex panel may add little forecasting value even when individual panel variables remain meaningful for communication, clinical judgement, or training planning.

This distinction is especially important when monitoring is intended to be individualised. Between-player differences in baseline scores, reporting behaviour, and physical responses can be substantial, and analyses of football performance have shown that monitoring should consider individual variability rather than rely solely on team averages \citep{OlivaLozano2021}. Evidence from women's football should likewise be interpreted in population-specific contexts, because physical match demands and psychological performance factors can vary substantially across female players and competitive environments \citep{Baptista2022,Pettersen2022}. Longitudinal datasets that combine repeated subjective load and wellness reports are valuable for addressing these questions. SoccerMon is notable because it provides openly available repeated subjective and objective monitoring data from two elite women's football teams \citep{Midoglu2024,SoccerMonZenodo}.

Prior work associated with the SoccerMon data stream has shown that readiness trends can be modelled with time-series approaches \citep{Wiik2019,Kulakou2022}. Yet feasibility demonstrations and comparisons of algorithms do not establish whether a full daily monitoring panel improves forecasting beyond a personalised readiness-history baseline, nor whether such value transports from one team to another. The distinction matters because the goal of the present study is prediction rather than causal explanation \citep{Shmueli2010}.

Prediction studies should also be evaluated beyond a single randomly split test set. When the intended use is future forecasting, evaluation should preserve temporal ordering and account for correlated player- and team-level observations \citep{Bergmeir2018,Roberts2017,Debray2023TRIPODCluster}. Calibration is also essential because a model with apparently acceptable error can still produce predictions with inappropriate spread or systematic bias \citep{VanCalster2019,Collins2024}. Cross-team testing further probes whether a model learned from one monitoring environment is transportable to another rather than merely exploiting stable team-specific patterns; performance in a new setting may require updating or recalibration \citep{Janssen2008}.

Accordingly, the primary aim of this study was to assess whether a daily workload--wellness panel adds incremental value beyond a personalised readiness-history baseline for forecasting next-day self-reported readiness in elite women's football. The secondary aim was to evaluate the stability of this incremental-value comparison under within-team temporal validation and bidirectional two-team transportability testing. We hypothesised that any incremental value of the full panel would need to be demonstrated consistently across these settings to be considered practically robust.
